\chapter*{第 6 章\quad 相似理论与量纲分析·孔珑题选}
\addcontentsline{toc}{chapter}{第 6 章\quad 相似理论与量纲分析·孔珑题选}
\markboth{第 6 章\quad 相似理论与量纲分析·孔珑题选}{}

\begin{tishi}
\noindent\textbf{题源与归章}\quad 本章题目取自孔珑《工程流体力学》（第四版）配套《工程流体力学 学习指导及习题解答》第 5 章“相似原理和量纲分析”5.3 节课后习题，原书题号 5-1 至 5-17，共 17 题，全部归属西华 818 参考书赵琴《流体力学与流体机械》第 6 章“相似理论与量纲分析”，一题不落、整章收录。

\noindent\textbf{本章星级口径}\quad 真题中本章内容只以小题出现：量纲分析与 $\pi$ 定理见于 E4 2016 填空第 3 题、2018 选择第 5 题、2019 选择第 5 题、2020 选择第 7 题以及 2020 判断第 10 题；$\Fr$ 的物理意义见于 2019 判断第 2 题；瑞利法见于 2019 判断第 8 题；可压缩流相似所用柯西数见于 2020 判断第 9 题；动力相似用雷诺准则见于 2014 选择第 10 题。E1 第 33 页只写“量纲转化、相似理论”，E2 明写“第六到第八章只有少量填空、判断、选择，掌握基本概念＋公式套用即可”，都没有把本章列为考试计算题重点，故本章不给五星：量纲分析与 $\pi$ 定理类记四星，相似准则与比尺换算类记三星。

\noindent\textbf{收录范围补充说明}\quad 5-2、5-15 的物理外形是堰，5-11 是风机、5-17 是水轮机；这几题考的都是相似比尺换算与 $\pi$ 定理，知识点属赵琴第 6 章，故予收录。风机、水轮机本身的结构与特性曲线不属本章要求。

\noindent\textbf{做题要求}\quad 比尺统一记作 $k$（如 $k_l$、$k_u$、$k_p$、$k_F$），原型量不加撇号、模型量加撇号；模型试验题先判断哪一种力起主要作用，再选相似准则；量纲分析题必须写出量纲方程并核对量纲一致性。
\end{tishi}

\begin{ti}\sloppy\emergencystretch=2em

\item \textbf{孔珑 5-1}\quad 试导出用基本量纲 L、T、M 表示的体积流量 $q_V$、质量流量 $q_m$、角速度 $\omega$、力矩 $M$、功 $W$、功率 $P$ 的量纲。
\xstarfull{4}\yiju{E4 2020 判断第 10 题考量纲、E4 2016 填空第 3 题考量纲分析；E5 题库选择题考运动黏滞系数量纲；E2 明写本章“掌握基本概念＋公式套用”}\nandu{易}

\item \textbf{孔珑 5-2}\quad 用模型研究溢流堰的流动，采用长度比例尺 $k_l=1/20$。
\tuyi{〔原书图 5-8，本册不重绘〕溢流堰纵剖面示意图：堰顶为一水平直线，上游水位高于堰顶，堰上水头 $h$ 自堰顶量至上游自由水面；水流越过堰顶后向下游跌落。模型与原型为几何相似，长度比例尺 $k_l=1/20$。}
\begin{xiaowen}
\item 已知原型堰上水头 $h=3\ \mathrm{m}$，试求模型的堰上水头 $h'$；
\item 测得模型上的体积流量 $q'_V=0.19\ \mathrm{m^3/s}$，试求原型上的体积流量 $q_V$；
\item 测得模型堰顶的计示压强 $p'_e=-1960\ \mathrm{Pa}$，试求原型堰顶的计示压强 $p_e$。
\end{xiaowen}
\xstarfull{3}\yiju{E4 2019 判断第 2 题考弗劳德数的物理意义（重力起主要作用的相似准则）；E2 明写本章“掌握基本概念＋公式套用”}\nandu{中}

\item \textbf{孔珑 5-3}\quad 有一内径 $d=200\ \mathrm{mm}$ 的圆管，输送运动黏度 $\nu=4.0\times10^{-5}\ \mathrm{m^2/s}$ 的油，其体积流量 $q_V=0.12\ \mathrm{m^3/s}$。若用内径 $d'=50\ \mathrm{mm}$ 的圆管并分别用 $20^{\circ}\mathrm{C}$ 的水和 $20^{\circ}\mathrm{C}$ 的空气做模型试验，试求流动相似时模型管内应有的体积流量。
（提示：$20^{\circ}\mathrm{C}$ 水的运动黏度 $\nu'=1.007\times10^{-6}\ \mathrm{m^2/s}$，$20^{\circ}\mathrm{C}$ 空气的运动黏度 $\nu'=15.00\times10^{-6}\ \mathrm{m^2/s}$。）
\xstarfull{3}\yiju{E4 2014 选择第 10 题考有压管流动力相似用雷诺准则；E5 题库选择题考雷诺准则；E2 明写本章“掌握基本概念＋公式套用”}\nandu{中}

\item \textbf{孔珑 5-4}\quad 将一高层建筑物的几何相似模型放在开口风洞中吹风，风速为 $u'=10\ \mathrm{m/s}$，测得模型迎风面点 1 处的计示压强 $p'_{1e}=980\ \mathrm{Pa}$，背风面点 2 处的计示压强 $p'_{2e}=-49\ \mathrm{Pa}$。试求该建筑物在 $u=30\ \mathrm{m/s}$ 的强风作用下对应点的计示压强。
\xstarfull{3}\yiju{E2 明写本章“掌握基本概念＋公式套用”；压强相似属欧拉准则，与 E4 2014 选择第 10 题的动力相似同一考法}\nandu{易}

\item \textbf{孔珑 5-5}\quad 长度比例尺 $k_l=1/40$ 的船模，当牵引速度 $u'=0.54\ \mathrm{m/s}$ 时，测得波阻 $F'_w=1.1\ \mathrm{N}$。如不计黏性影响，试求原型船的速度、波阻及消耗的功率。
\xstarfull{3}\yiju{E4 2019 判断第 2 题考弗劳德数的物理意义；E2 明写本章“掌握基本概念＋公式套用”}\nandu{中}

\item \textbf{孔珑 5-6}\quad 长度比例尺 $k_l=1/225$ 的模型水库，开闸后完全放空库水的时间是 $t'=4\ \mathrm{min}$。试求原型水库放空库水的时间 $t$。
\xstarfull{3}\yiju{E4 2019 判断第 2 题考弗劳德数的物理意义（重力相似下时间比尺）；E2 明写本章“掌握基本概念＋公式套用”}\nandu{易}

\item \textbf{孔珑 5-7}\quad 新设计的汽车高 $1.5\ \mathrm{m}$，最大行驶速度为 $108\ \mathrm{km/h}$，拟在风洞中进行模型试验。已知风洞试验段的最大风速为 $45\ \mathrm{m/s}$，试求模型的高度。在该风速下测得模型的风阻力为 $1500\ \mathrm{N}$，试求原型在最大行驶速度时的风阻。
（提示：汽车与风洞中的模型都在空气中，且温度相同，故 $\nu'=\nu$。）
\xstarfull{3}\yiju{E4 2014 选择第 10 题考动力相似用雷诺准则；E5 题库选择题考雷诺准则；E2 明写本章“掌握基本概念＋公式套用”}\nandu{中}

\item \textbf{孔珑 5-8}\quad 在管道内以 $u=20\ \mathrm{m/s}$ 的速度输送密度 $\rho=1.86\ \mathrm{kg/m^3}$、运动黏度 $\nu=1.3\times10^{-5}\ \mathrm{m^2/s}$ 的天然气，为了预测沿管道的压强降，采用水模型试验。取长度比例尺 $k_l=1/10$，已知水的密度 $\rho_w=998\ \mathrm{kg/m^3}$、运动黏度 $\nu_w=1.007\times10^{-6}\ \mathrm{m^2/s}$。为保证流动相似，模型内水的流速应等于多少？已测得模型每 $0.1\ \mathrm{m}$ 管长的压强降 $\Delta p'=1000\ \mathrm{Pa}$，天然气管道每米的压强降等于多少？
\xstarfull{3}\yiju{E4 2014 选择第 10 题考有压管流动力相似用雷诺准则；E5 题库选择题考雷诺准则；E2 明写本章“掌握基本概念＋公式套用”}\nandu{难}

\item \textbf{孔珑 5-9}\quad 烟气在 $600^{\circ}\mathrm{C}$ 的热处理炉中的流动情况拟用水模型来进行研究。已知烟气的运动黏度 $\nu=9.0\times10^{-5}\ \mathrm{m^2/s}$，长度比例尺 $k_l=1/10$，水温为 $10^{\circ}\mathrm{C}$，试求速度比例尺。
（提示：$10^{\circ}\mathrm{C}$ 水的运动黏度 $\nu'=1.308\times10^{-6}\ \mathrm{m^2/s}$。）
\xstarfull{3}\yiju{E4 2014 选择第 10 题考动力相似用雷诺准则；E2 明写本章“掌握基本概念＋公式套用”}\nandu{易}

\item \textbf{孔珑 5-10}\quad 某飞机的机翼弦长 $b=1500\ \mathrm{mm}$，在气压 $p_a=10^5\ \mathrm{Pa}$、气温 $t=10^{\circ}\mathrm{C}$ 的大气中以 $u=180\ \mathrm{km/h}$ 的速度飞行，拟在风洞中用模型试验测定翼型阻力，采用长度比例尺 $k_l=1/3$。
（提示：由教材图 2-5 查得空气在 $10^{\circ}\mathrm{C}$ 时的动力黏度 $\mu=1.75\times10^{-5}\ \mathrm{Pa\cdot s}$，在 $25^{\circ}\mathrm{C}$ 时的动力黏度 $\mu'=1.82\times10^{-5}\ \mathrm{Pa\cdot s}$。）
\begin{xiaowen}
\item 如果用开口风洞，已知试验段的气压 $p'=101325\ \mathrm{Pa}$、气温 $t'=25^{\circ}\mathrm{C}$，试验段的风速应等于多少？这样的试验有什么问题？
\item 如果用压力风洞，试验段的气压 $p''=1\ \mathrm{MPa}$、气温 $t''=30^{\circ}\mathrm{C}$、动力黏度 $\mu''=1.854\times10^{-5}\ \mathrm{Pa\cdot s}$，试验段的风速应等于多少？
\end{xiaowen}
\xstarfull{3}\yiju{E4 2020 判断第 9 题考可压缩流的柯西数（可压缩性相似）；E2 明写本章“掌握基本概念＋公式套用”}\nandu{难}

\item \textbf{孔珑 5-11}\quad 低压轴流风机的叶轮直径 $d=0.4\ \mathrm{m}$，转速 $n=1400\ \mathrm{r/min}$，体积流量 $q_V=1.39\ \mathrm{m^3/s}$，全压 $p=128\ \mathrm{Pa}$，效率 $\eta=70\%$，空气密度 $\rho=1.20\ \mathrm{kg/m^3}$，问消耗的功率 $P$ 等于多少？在保证流动相似和假定风机效率不变的情况下，试确定作下列三种变动情况下的 $q'_V$、$p'$ 和 $P'$ 值：
\begin{xiaowen}
\item 转速 $n$ 变为 $2800\ \mathrm{r/min}$；
\item 风机相似放大，叶轮直径 $d'$ 变为 $0.8\ \mathrm{m}$；
\item 空气密度 $\rho'$ 变为 $1.29\ \mathrm{kg/m^3}$。
\end{xiaowen}
\xstarfull{3}\yiju{E5 题库问答题考泵与风机的标准条件及相似换算；E2 明写本章“掌握基本概念＋公式套用”}\nandu{中}

\item \textbf{孔珑 5-12}\quad 流体通过水平毛细管的体积流量 $q_V$ 与管径 $d$、动力黏度 $\mu$、压强梯度 $\Delta p/l$ 有关，试导出流量的表达式。
\xstarfull{3}\yiju{E4 2019 判断第 8 题考瑞利法待求量纲指数不超过 4 个；E2 明写本章“掌握基本概念＋公式套用”}\nandu{中}

\item \textbf{孔珑 5-13}\quad 薄壁孔口出流的流速 $v$ 与孔口直径 $d$、孔口上水头 $H$、流体密度 $\rho$、动力黏度 $\mu$、表面张力 $\sigma$、重力加速度 $g$ 有关，试导出孔口出流速度的表达式。
\xstarfull{4}\yiju{E4 2016 填空第 3 题、2018 选择第 5 题、2019 选择第 5 题、2020 选择第 7 题均考量纲分析与 $\pi$ 定理（但都是小题，故按本册星级表记四星）}\nandu{难}

\item \textbf{孔珑 5-14}\quad 小球在不可压缩黏性流体中运动的阻力 $F_D$ 与小球直径 $d$、等速运动的速度 $v$、流体的密度 $\rho$、动力黏度 $\mu$ 有关，试导出阻力的表达式。
\xstarfull{4}\yiju{E4 2016 填空第 3 题、2018 选择第 5 题、2019 选择第 5 题、2020 选择第 7 题均考量纲分析与 $\pi$ 定理；本题即布金汉定理的直接套用}\nandu{中}

\item \textbf{孔珑 5-15}\quad 通过三角形堰的体积流量 $q_V$ 与堰上水头 $H$、槽口的半顶角 $\theta$、重力加速度 $g$、液体的密度 $\rho$、动力黏度 $\mu$、表面张力 $\sigma$ 有关，试导出流量的表达式。
\tuyi{〔原书图 5-9，本册不重绘〕三角形堰（V 形缺口堰）正视图：堰板上开一个倒三角形缺口，缺口两斜边的夹角为 $2\theta$，半顶角为 $\theta$；堰上水头 $H$ 自缺口的顶点量至上游自由水面；水经三角形缺口溢流。}
\xstarfull{4}\yiju{E4 2016 填空第 3 题、2018 选择第 5 题、2019 选择第 5 题、2020 选择第 7 题均考量纲分析与 $\pi$ 定理；E2 明写本章“掌握基本概念＋公式套用”}\nandu{难}

\item \textbf{孔珑 5-16}\quad 流体通过孔板流量计的体积流量 $q_V$ 与孔板前后的压差 $\Delta p$、管道的内径 $d$、管内流速 $u$、孔板的孔径 $d_0$、流体的密度 $\rho$、动力黏度 $\mu$ 有关，试导出流量 $q_V$ 的表达式。
\xstarfull{4}\yiju{E4 2016 填空第 3 题、2018 选择第 5 题、2019 选择第 5 题、2020 选择第 7 题均考量纲分析与 $\pi$ 定理；E2 明写本章“掌握基本概念＋公式套用”}\nandu{难}

\item \textbf{孔珑 5-17}\quad 水轮机的功率 $P$ 与叶轮直径 $d$、叶片宽度 $b$、转速 $n$、有效水头 $H$、水的密度 $\rho$、动力黏度 $\mu$、重力加速度 $g$ 有关，试导出水轮机功率的表达式。
\xstarfull{4}\yiju{E4 2016 填空第 3 题、2018 选择第 5 题、2019 选择第 5 题、2020 选择第 7 题均考量纲分析与 $\pi$ 定理；本题为 $\pi$ 定理综合应用题}\nandu{难}

\end{ti}

\begin{tishi}
\noindent\textbf{本章是否有略去的题号}\quad 无。孔珑第 5 章 5.3 节课后习题共 5-1 至 5-17 共 17 题，题题对应赵琴第 6 章“相似理论与量纲分析”，全部收录，未作删减。

\noindent\textbf{题面 OCR 校订清单}\quad 本册题目来自影印本 OCR，已按物理意义订正：① 5-1 中原印“$\dim q_V=\mathrm{L^2T^{-1}}$”应为 $\mathrm{L^3T^{-1}}$，原印“$\dim M=\mathrm{MLT^{-1}}$”应为 $\mathrm{ML^2T^{-2}}$（力矩与功同量纲）；② 影印本中体积流量单位多被识别成 $\mathrm{m^2/s}$，一律按 $\mathrm{m^3/s}$ 订正；③ 密度单位多被识别成 $\mathrm{kg/m^2}$，一律按 $\mathrm{kg/m^3}$ 订正；④ 5-10 中原印“气压 $p_a=10\ \mathrm{Pa}$”应为 $10^5\ \mathrm{Pa}$；⑤ 5-5 中原印功率“$2.404\times10\ \mathrm{W}$”应为 $2.404\times10^5\ \mathrm{W}$；⑥ 希腊字母与比尺下标极易误识：$\lambda_l$ 应为比尺 $k_l$（本册比尺统一写作 $k$），$\pi$ 常被识别成 n，$\rho$ 常被识别成 p，$\nu$ 常被识别成 v 或 u，$\mu$ 常被识别成 u；⑦ 5-11 中原印“效率 $\eta=70\%$”曾被识别成 n；⑧ 5-11 的第 3 种变动中原印“变为 1.29”指的是空气密度 $\rho'$ 变为 $1.29\ \mathrm{kg/m^3}$。

\noindent\textbf{题面残缺说明}\quad 5-2 与 5-15 的原书插图（图 5-8、图 5-9）为影印图片，本册不重绘，已用图示命令作文字描述（给出图形含义与已知量），题目所给数据完整，不影响求解。各题题面数据齐全，无残缺。

\noindent\textbf{超出考纲的说明}\quad 5-11 的三种变动（变转速、几何放大、变密度）在工程上即风机的比例定律与相似换算，题库（E5 问答题）考过“泵与风机的标准条件及相似换算”，故保留；但风机、水轮机的效率曲线、特性曲线、比转速等内容属赵琴第 11 至 14 章（复试范围），本章不作要求，只需会算比尺。
\end{tishi}
